
\documentclass{article}
\usepackage{colortbl}
\usepackage{makecell}
\usepackage{multirow}
\usepackage{supertabular}

\begin{document}

\newcounter{utterance}

\centering \large Interaction Transcript for game `privateshared', experiment `things{-}places', episode 1 with Qwen3{-}30B{-}A3B{-}t0.0.
\vspace{24pt}

{ \footnotesize  \setcounter{utterance}{1}
\setlength{\tabcolsep}{0pt}
\begin{supertabular}{c@{$\;$}|p{.15\linewidth}@{}p{.15\linewidth}p{.15\linewidth}p{.15\linewidth}p{.15\linewidth}p{.15\linewidth}}
    \# & $\;$A & \multicolumn{4}{c}{Game Master} & $\;\:$B\\
    \hline

    \theutterance \stepcounter{utterance}  
    & & \multicolumn{4}{p{0.6\linewidth}}{
        \cellcolor[rgb]{0.9,0.9,0.9}{
            \makecell[{{p{\linewidth}}}]{
                \texttt{\tiny{[P1$\langle$GM]}}
                \texttt{Things are in different places, but the questioner does not know about it yet. Here is the situation:} \\
\\ 
\texttt{WHAT: Things at places} \\
\texttt{LEFT: books} \\
\texttt{RIGHT: strawberry} \\
\texttt{TOP: boron} \\
\texttt{BOTTOM: orange} \\
\texttt{CENTER: plate} \\
\texttt{NORTHWEST: eraser} \\
\texttt{NORTHEAST: horse} \\
\texttt{SOUTHWEST: evening gown} \\
\texttt{SOUTHEAST: copper} \\
\texttt{HERE: primer} \\
\texttt{THERE: blueberry} \\
\texttt{NOWHERE: ukulele} \\
\texttt{EVERYWHERE: xylophone} \\
\texttt{INSIDE: mango} \\
\texttt{OUTSIDE: daffodil} \\
\\ 
\texttt{Questions from the questioner will start with QUESTIONER. Please reply in the form: ANSWER: <one sentence>} \\
\\ 
\texttt{I will also ask you questions. These questions will start with ME: . Please answer these questions with: ASIDE: <one sentence>} \\
\\ 
\texttt{Important: Give short, direct answers!} \\
\\ 
\texttt{Let us start.} \\
            }
        }
    }
    & \\ \\

    \theutterance \stepcounter{utterance}  
    & \multicolumn{4}{p{0.6\linewidth}}{
        \cellcolor[rgb]{0.9,0.9,0.9}{
            \makecell[{{p{\linewidth}}}]{
                \texttt{\tiny{[P1$\rangle$GM]}}
                \texttt{Ok.} \\
            }
        }
    }
    & & \\ \\

    \theutterance \stepcounter{utterance}  
    & & \multicolumn{4}{p{0.6\linewidth}}{
        \cellcolor[rgb]{0.9,0.9,0.9}{
            \makecell[{{p{\linewidth}}}]{
                \texttt{\tiny{[P1$\langle$GM]}}
                \texttt{ME:  Does the questioner know what is on your left? Answer this question with yes or no and then stop.} \\
            }
        }
    }
    & \\ \\

    \theutterance \stepcounter{utterance}  
    & \multicolumn{4}{p{0.6\linewidth}}{
        \cellcolor[rgb]{0.9,0.9,0.9}{
            \makecell[{{p{\linewidth}}}]{
                \texttt{\tiny{[P1$\rangle$GM]}}
                \texttt{<think>} \\
\texttt{Okay, the user is asking if the questioner knows what's on my left. Let me check the information provided.} \\
\\ 
\texttt{In the setup, the left place is listed as "books". The questioner's knowledge isn't explicitly stated. The key point here is whether the questioner has been informed about the items in each location. Since the questioner hasn't asked anything yet and the initial instruction says "the questioner does not know about it yet," I should assume they aren't aware. So the answer is no. I need to respond with "ANSWER: no" and stop after that.} \\
\texttt{</think>} \\
\\ 
\texttt{ANSWER: no} \\
            }
        }
    }
    & & \\ \\

    \theutterance \stepcounter{utterance}  
    & & & \multicolumn{2}{p{0.3\linewidth}}{
        \cellcolor[rgb]{0.9,0.9,0.9}{
            \makecell[{{p{\linewidth}}}]{
                \texttt{\tiny{[GM$|$GM]}}
                \texttt{NA} \\
            }
        }
    }
    & & \\ \\

    \theutterance \stepcounter{utterance}  
    & & \multicolumn{4}{p{0.6\linewidth}}{
        \cellcolor[rgb]{0.9,0.9,0.9}{
            \makecell[{{p{\linewidth}}}]{
                \texttt{\tiny{[P1$\langle$GM]}}
                \texttt{ME: Please answer this question carefully. Does the questioner know what is on your left? Answer this question with yes or no and then stop.} \\
            }
        }
    }
    & \\ \\

    \theutterance \stepcounter{utterance}  
    & \multicolumn{4}{p{0.6\linewidth}}{
        \cellcolor[rgb]{0.9,0.9,0.9}{
            \makecell[{{p{\linewidth}}}]{
                \texttt{\tiny{[P1$\rangle$GM]}}
                \texttt{<think>} \\
\texttt{Okay, the user is asking if the questioner knows what's on my left. Let me check the setup again.} \\
\\ 
\texttt{The given information lists "LEFT: books". So, according to the data provided, the left side has books. But the question is about the questioner's knowledge. The initial statement says, "Things are in different places, but the questioner does not know about it yet." That implies the questioner isn't aware of the current placements. } \\
\\ 
\texttt{So even though the left is books, the questioner doesn't know that. The answer should be no. The user also specified to answer with yes or no and then stop. So the response should be "no" followed by a period. I need to make sure I don't add anything else. Let me confirm once more. The key part is "the questioner does not know about it yet." Yep, that's clear. So the answer is no.} \\
\texttt{</think>} \\
\\ 
\texttt{ANSWER: no.} \\
            }
        }
    }
    & & \\ \\

    \theutterance \stepcounter{utterance}  
    & & & \multicolumn{2}{p{0.3\linewidth}}{
        \cellcolor[rgb]{0.9,0.9,0.9}{
            \makecell[{{p{\linewidth}}}]{
                \texttt{\tiny{[GM$|$GM]}}
                \texttt{NA} \\
            }
        }
    }
    & & \\ \\

    \theutterance \stepcounter{utterance}  
    & & \multicolumn{4}{p{0.6\linewidth}}{
        \cellcolor[rgb]{0.9,0.9,0.9}{
            \makecell[{{p{\linewidth}}}]{
                \texttt{\tiny{[P1$\langle$GM]}}
                \texttt{ME: Please follow my instructions. Does the questioner know what is on your left? Answer this question with yes or no and then stop.} \\
            }
        }
    }
    & \\ \\

    \theutterance \stepcounter{utterance}  
    & \multicolumn{4}{p{0.6\linewidth}}{
        \cellcolor[rgb]{0.9,0.9,0.9}{
            \makecell[{{p{\linewidth}}}]{
                \texttt{\tiny{[P1$\rangle$GM]}}
                \texttt{<think>} \\
\texttt{Okay, the user is asking if the questioner knows what's on my left. Let me check the setup again.} \\
\\ 
\texttt{The given information lists "LEFT: books". So according to the data provided, the left side has books. But the question says that the questioner doesn't know about the things being in different places yet. So even though the answer is "books", the questioner isn't aware of that. } \\
\\ 
\texttt{The user wants a yes or no answer followed by stopping. Since the questioner doesn't know, the answer should be no. I need to make sure I don't add any extra info and just respond with "no" as the answer.} \\
\texttt{</think>} \\
\\ 
\texttt{ANSWER: no} \\
            }
        }
    }
    & & \\ \\

    \theutterance \stepcounter{utterance}  
    & & & \multicolumn{2}{p{0.3\linewidth}}{
        \cellcolor[rgb]{0.9,0.9,0.9}{
            \makecell[{{p{\linewidth}}}]{
                \texttt{\tiny{[GM$|$GM]}}
                \texttt{NA} \\
            }
        }
    }
    & & \\ \\

    \theutterance \stepcounter{utterance}  
    & & \multicolumn{4}{p{0.6\linewidth}}{
        \cellcolor[rgb]{0.9,0.9,0.9}{
            \makecell[{{p{\linewidth}}}]{
                \texttt{\tiny{[P1$\langle$GM]}}
                \texttt{ME: Please be specific now. Does the questioner know what is on your left? Answer this question with yes or no and then stop.} \\
            }
        }
    }
    & \\ \\

    \theutterance \stepcounter{utterance}  
    & \multicolumn{4}{p{0.6\linewidth}}{
        \cellcolor[rgb]{0.9,0.9,0.9}{
            \makecell[{{p{\linewidth}}}]{
                \texttt{\tiny{[P1$\rangle$GM]}}
                \texttt{<think>} \\
\texttt{Okay, the user is asking if the questioner knows what's on my left. Let me check the setup again.} \\
\\ 
\texttt{The given information lists "LEFT: books". So according to the data, the left side has books. But the question says the questioner doesn't know about the things yet. So even though the answer is books, the questioner isn't aware. The user wants a yes or no answer followed by stopping. So the answer should be no because the questioner hasn't been informed yet.} \\
\texttt{</think>} \\
\\ 
\texttt{ANSWER: No.} \\
            }
        }
    }
    & & \\ \\

    \theutterance \stepcounter{utterance}  
    & & & \multicolumn{2}{p{0.3\linewidth}}{
        \cellcolor[rgb]{0.9,0.9,0.9}{
            \makecell[{{p{\linewidth}}}]{
                \texttt{\tiny{[GM$|$GM]}}
                \texttt{NA} \\
            }
        }
    }
    & & \\ \\

    \theutterance \stepcounter{utterance}  
    & & \multicolumn{4}{p{0.6\linewidth}}{
        \cellcolor[rgb]{0.9,0.9,0.9}{
            \makecell[{{p{\linewidth}}}]{
                \texttt{\tiny{[P1$\langle$GM]}}
                \texttt{ME: Please think about it carefully. Does the questioner know what is on your left? Answer this question with yes or no and then stop.} \\
            }
        }
    }
    & \\ \\

    \theutterance \stepcounter{utterance}  
    & \multicolumn{4}{p{0.6\linewidth}}{
        \cellcolor[rgb]{0.9,0.9,0.9}{
            \makecell[{{p{\linewidth}}}]{
                \texttt{\tiny{[P1$\rangle$GM]}}
                \texttt{<think>} \\
\texttt{Okay, the user is asking if the questioner knows what's on my left. Let me check the information provided.} \\
\\ 
\texttt{In the setup, the user listed various items in different positions. The LEFT is books. But the questioner hasn't been told anything yet. The initial message says the questioner doesn't know the locations. So the questioner isn't aware of the items in each place. } \\
\\ 
\texttt{The user wants a yes or no answer. Since the questioner hasn't been informed, the answer should be no. I need to make sure I'm not adding any extra info. Just a straightforward yes or no. The answer is no.} \\
\texttt{</think>} \\
\\ 
\texttt{ANSWER: No.} \\
            }
        }
    }
    & & \\ \\

    \theutterance \stepcounter{utterance}  
    & & & \multicolumn{2}{p{0.3\linewidth}}{
        \cellcolor[rgb]{0.9,0.9,0.9}{
            \makecell[{{p{\linewidth}}}]{
                \texttt{\tiny{[GM$|$GM]}}
                \texttt{NA} \\
            }
        }
    }
    & & \\ \\

    \theutterance \stepcounter{utterance}  
    & & & \multicolumn{2}{p{0.3\linewidth}}{
        \cellcolor[rgb]{0.9,0.9,0.9}{
            \makecell[{{p{\linewidth}}}]{
                \texttt{\tiny{[GM$|$GM]}}
                \texttt{Answer for LEFT invalid after max attempts.} \\
            }
        }
    }
    & & \\ \\

    \theutterance \stepcounter{utterance}  
    & & & \multicolumn{2}{p{0.3\linewidth}}{
        \cellcolor[rgb]{0.9,0.9,0.9}{
            \makecell[{{p{\linewidth}}}]{
                \texttt{\tiny{[GM$|$GM]}}
                \texttt{Answer is incorrect.} \\
            }
        }
    }
    & & \\ \\

    \theutterance \stepcounter{utterance}  
    & & & \multicolumn{2}{p{0.3\linewidth}}{
        \cellcolor[rgb]{0.9,0.9,0.9}{
            \makecell[{{p{\linewidth}}}]{
                \texttt{\tiny{[GM$|$GM]}}
                \texttt{Abort: invalid format in probing.} \\
            }
        }
    }
    & & \\ \\

    \theutterance \stepcounter{utterance}  
    & & \multicolumn{4}{p{0.6\linewidth}}{
        \cellcolor[rgb]{0.9,0.9,0.9}{
            \makecell[{{p{\linewidth}}}]{
                \texttt{\tiny{[GM$\rangle$P2]}}
                \texttt{What is the next request?} \\
            }
        }
    }
    & \\ \\

    \theutterance \stepcounter{utterance}  
    & & & \multicolumn{4}{p{0.6\linewidth}}{
        \cellcolor[rgb]{0.9,0.9,0.9}{
            \makecell[{{p{\linewidth}}}]{
                \texttt{\tiny{[GM$\langle$P2]}}
                \texttt{What is at the center?} \\
            }
        }
    }
     \\ \\

    \theutterance \stepcounter{utterance}  
    & & & \multicolumn{2}{p{0.3\linewidth}}{
        \cellcolor[rgb]{0.9,0.9,0.9}{
            \makecell[{{p{\linewidth}}}]{
                \texttt{\tiny{[GM$|$GM]}}
                \texttt{Game finished.} \\
            }
        }
    }
    & & \\ \\

\end{supertabular}
}

\end{document}
